% !PRINTED 418
\weberpage{418}
%<<CONT
\gesperrt{ganze rationale Functionen von $x$ mit rationalen Zahlenco\"efficienten sind.}

Dieser wichtige Satz ist für den Fall, dass $n$ eine Primzahl ist, von
\gesperrt{Gauss} aufgestellt und zuerst bewiesen. Später sind noch viele
andere Beweise gegeben worden, die zum Theil allgemeiner und einfacher sind.
Wir folgen hier einem Beweis von \gesperrt{Kronecker}.

In der Function $X_n$ hat, wie wir vorhin gesehen haben, die höchste Potenz
von $x$ den Co\"efficienten $1$, während die übrigen Co\"efficienten ganze
Zahlen sind. Nehmen wir nun an, es könne $X_n$ in zwei Factoren mit
rationalen Co\"efficienten

\weq{1}{X_n = \varphi(x)\,\psi(x)}

zerlegt werden, so muss das Product der Co\"efficienten der höchsten Potenzen
von $x$ in $\varphi$ und $\psi$ gleich $1$ sein, und wir können also
annehmen, dass diese höchsten Co\"efficienten in $\varphi$ und $\psi$ selbst
den Werth $1$ haben (da wir anderenfalls durch das Product der beiden
Co\"efficienten dividiren würden). Dann folgt aber aus dem
\gesperrt{Gauss}'schen Satz (\S.\,2), dass die übrigen Co\"efficienten von
$\varphi(x)$ und $\psi(x)$ \gesperrt{ganze Zahlen} sein müssen, dass also $\varphi(x)$
und $\psi(x)$, wenn für $x$ eine ganze Zahl gesetzt wird, selbst in ganze
Zahlen übergehen müssen.

Setzen wir also $x = 1$, so folgt:
\[
\varphi(1)\,\psi(1) = X_n(1),
\]
also nach Satz V. gleich $1$ oder gleich einer Primzahl $p$. Demnach muss von
den beiden Factoren $\varphi(1)$, $\psi(1)$ gewiss einer den Werth $\pm 1$
haben, während der andere gleich $\pm 1$ oder $\pm p$ ist. Wir nehmen also
an, es sei

\weq{2}{\varphi(1) = \pm 1.}

Die Wurzeln $r$ der Gleichung $X_n = 0$ vertheilen sich nun auf die beiden
Gleichungen $\varphi = 0$, $\psi = 0$, und da keiner der beiden Factoren
$\varphi$, $\psi$ von $x$ unabhängig ist, so muss es unter den primitiven
$n^{\mathrm{ten}}$ Einheitswurzeln $r$ gewiss eine geben, die wir mit $\rho$
bezeichnen wollen, für die
\[
\varphi(\rho) = 0
\]
ist.

Wenn nun $r$ irgend eine primitive $n^{\mathrm{te}}$ Einheitswurzel
bedeutet, so sind unter den Potenzen
\[
r,\ r^2,\ r^3\ \dots\ r^{n-1}
\]
alle $n^{\mathrm{ten}}$ Einheitswurzeln mit Ausnahme von $1$, also auch
$\rho$ enthalten, und es folgt also, dass das Product
%CONT>>
